% Chapter 9: Importance sampling.
\providecommand{\Qt}{\widetilde{\mathbb Q}}
\providecommand{\Et}{\widetilde{\mathbb E}}

A Monte Carlo estimator spends its samples where the pricing measure puts its mass. For an
at-the-money call that is roughly the right place. For a call struck three standard
deviations out of the money, more than $99.8\%$ of the paths pay nothing, and the price is
decided by the few that do. \emph{Importance sampling} draws the samples from another law,
chosen to put them where the payoff lives, and removes the bias by weighting each sample
with the likelihood ratio of the two laws. The same ratio drives the likelihood-ratio Greeks
of \cref{ch:13-mc-greeks}; here it is used the other way round.

The method can cut the variance by orders of magnitude, and it can also make it infinite
while the estimator still looks well behaved; every gain below therefore comes with its
failure mode. The code is \coderef{code/src/mfe230d/importance.py}.

Throughout, the test case is $S=K=100$, $T=1$, $r=1\%$, $y=2\%$, $\sigma=20\%$, with
$\log(S_T/S)\sim\mathcal N(\nu T,\sigma^2T)$ under $\Q$, $\nu=r-y-\sigma^2/2=-0.03$, and
call price $7.36429$. In standard units $S_T=S\exp(\nu T+sZ)$ with $s=\sigma\sqrt T$ and
$Z\sim\mathcal N(0,1)$, and the call pays exactly when $Z>d$, where
\[
  d=\frac{\log(K/S)-\nu T}{\sigma\sqrt T}=-d_-
\]
is the distance from the centre of the law of $Z$ to the strike, in standard deviations.

\section{The likelihood ratio}\label{sec:is-lr}

Let $X$ be a random vector in $\R^k$ with density $q$ under $\Q$, and let $f:\R^k\to\R$ be
Borel measurable with $\EQ|f(X)|<\infty$; we want $\mu=\EQ[f(X)]$. In pricing, $X$ is the
terminal log-price or the vector of normals that drives a simulated path, and $f$ includes
the discount factor. Let $\tilde q$ be another probability density on $\R^k$, the
\emph{sampling density}, and write $\Qt$ and $\Et$ for probabilities and expectations when
$X$ has density $\tilde q$. The \emph{likelihood ratio}, or \emph{weight}, is the Borel
function
\[
  L(x)=\frac{q(x)}{\tilde q(x)}\ \text{ if }\tilde q(x)>0,\qquad L(x)=0\ \text{ otherwise}.
\]
Everything below holds verbatim for discrete laws, with densities with respect to counting
measure (probability mass functions).

\begin{proposition}[The importance-sampling estimator]\label{prop:is-likelihood-ratio}
Assume the \emph{support condition}
\begin{equation}\label{eq:is-support}
  \tilde q(x)>0\quad\text{for every }x\text{ with }f(x)q(x)\neq0 .
\end{equation}
\begin{enumerate}[label=(\roman*)]
  \item $f(X)L(X)$ is $\Qt$-integrable and $\Et[f(X)L(X)]=\EQ[f(X)]=\mu$.
  \item If $X_1,X_2,\dots$ are independent with density $\tilde q$, the
    \emph{importance-sampling estimator}
    $\hat\mu^{\mathrm{is}}_n=\frac1n\sum_{i=1}^nf(X_i)L(X_i)$ is unbiased and converges to
    $\mu$ almost surely.
  \item Its variance is $\sigma_{\mathrm{is}}^2/n$, where
    \begin{equation}\label{eq:is-variance}
      \sigma_{\mathrm{is}}^2+\mu^2=\Et\big[f(X)^2L(X)^2\big]=\EQ\big[f(X)^2L(X)\big]
      =\int_{\R^k}\frac{f(x)^2q(x)^2}{\tilde q(x)}\,\dd x\ \in[0,\infty],
    \end{equation}
    with the integral taken over $\{\tilde q>0\}$.
  \item Without \eqref{eq:is-support}, $\Et[f(X)L(X)]=\EQ\big[f(X)\1_{\{\tilde q(X)>0\}}\big]$:
    the estimator converges to the wrong number, and nothing in the samples reveals it.
\end{enumerate}
\end{proposition}

\begin{proof}
For any Borel $g\ge0$, or $g$ with $\EQ|g(X)|<\infty$,
\[
  \Et[g(X)L(X)]=\int_{\{\tilde q>0\}}g(x)\frac{q(x)}{\tilde q(x)}\,\tilde q(x)\,\dd x
  =\int_{\{\tilde q>0\}}g(x)q(x)\,\dd x=\EQ\big[g(X)\1_{\{\tilde q(X)>0\}}\big].
\]
With $g=|f|$ this shows integrability, and with $g=f$ it gives (iv). Under
\eqref{eq:is-support}, $f(x)q(x)=0$ on $\{\tilde q=0\}$, so the indicator can be dropped:
this is (i). (ii) follows from (i) and the strong law of large numbers for the i.i.d.\
integrable variables $f(X_i)L(X_i)$. For (iii), apply the display to $g=f^2L\ge0$ (Tonelli):
$\Et[f^2L^2]=\EQ[f^2L\,\1_{\{\tilde q>0\}}]=\EQ[f^2L]$, since $L=0$ off $\{\tilde q>0\}$;
both sides may be infinite. Subtract $\mu^2$.
\end{proof}

If $0<\sigma^2_{\mathrm{is}}<\infty$, \cref{thm:mc-clt} gives confidence intervals with the
sample standard deviation $\hat s_n$ of the weighted samples; if $\sigma^2_{\mathrm{is}}=0$,
every sample equals $\mu$; if $\sigma^2_{\mathrm{is}}=\infty$, the estimator still converges,
but $\hat s_n\to\infty$ almost surely and the intervals mean nothing
(\cref{rem:mc-heavy-tails}). The weight is always integrable, $\Et[L]=\Q(\tilde q(X)>0)\le1$
($g=1$ in the proof), but its variance can be infinite.

The middle expression in \eqref{eq:is-variance} is the key to everything that follows: the
second moment of the payoff under the target law, weighted by $L$. Where the sampling
density is much smaller than the target density, $L$ is large and the payoff is punished;
where it is larger, $L<1$ and the payoff is discounted. A good sampling density is large
where $|f|q$ is large. The extra cost of a sample is one evaluation of $L$, usually one
exponential, so the efficiency gain over plain Monte Carlo (\cref{def:mc-efficiency}) is
close to the variance ratio $\sigma^2/\sigma^2_{\mathrm{is}}$. Under $\Qt$, discounted prices
are not martingales: the sampling law is a computational device, not a pricing measure.

\begin{proposition}[The zero-variance law]\label{prop:is-zero-variance}
Assume $0<\EQ|f(X)|<\infty$. For every sampling density satisfying \eqref{eq:is-support},
\[
  \sigma^2_{\mathrm{is}}\ \ge\ \big(\EQ|f(X)|\big)^2-\mu^2 ,
\]
with equality if and only if $\tilde q=|f|q/\EQ|f(X)|$ almost everywhere. In particular, if
$f\ge0$, the density $\tilde q^*=fq/\mu$ gives $f(X)L(X)=\mu$ for $\Qt$-almost every sample,
and it is the only sampling density with zero variance.
\end{proposition}

\begin{proof}
By Cauchy--Schwarz and \cref{prop:is-likelihood-ratio}(i) applied to $|f|$,
$\Et[(fL)^2]\ge\big(\Et|fL|\big)^2=\big(\EQ|f|\big)^2$, which is the inequality. Equality
holds if and only if $|f(X)|L(X)$ is $\Qt$-almost surely equal to a constant $c$, and then
$c=\Et|fL|=\EQ|f|>0$. On $\{\tilde q>0,\ fq=0\}$ we would have $|f|L=0\ne c$, so this set
has $\Qt$-probability zero; on $\{fq\ne0\}$, which lies in $\{\tilde q>0\}$ by
\eqref{eq:is-support}, $|f|q/\tilde q=c$, that is $\tilde q=|f|q/c$. Hence
$\tilde q=|f|q/c$ almost everywhere; conversely this density gives $|f|L=c$. For $f\ge0$,
$\EQ|f|=\mu$ and the bound is $0$.
\end{proof}

The zero-variance density needs $\mu$, the number we are computing, so it is never used as
such. What it gives is the shape to aim for: \emph{sample in proportion to payoff times
density}. For a call it is the risk-neutral density of $S_T$ on $\{S_T>K\}$, multiplied by
$S_T-K$ (\cref{fig:is-densities}). A payoff that changes sign cannot be estimated with zero
variance by one sampling law; writing $f=f^+-f^-$ and estimating the two parts with two laws
removes the obstruction.

\begin{figure}[ht]
  \centering
  \includegraphics[width=0.85\textwidth]{ch09-densities}
  \caption{The course's experiment at $K=100$. (a) The risk-neutral density of $S_T$ and the
    sampling densities with drift shifts $\alpha=0.1,0.2,0.3$ ($\beta=1$); dashed, the
    zero-variance density of \cref{prop:is-zero-variance}. (b) The weighted payoff
    $e^{-rT}(S_T-K)^+L$, whose $\Qt$-mean is the price $7.36$ (dotted) in every case: the
    shift turns the unbounded payoff into a bounded hump, much flatter under the sampling law.}
  \label{fig:is-densities}
\end{figure}

\section{Shifting and scaling the Gaussian}\label{sec:is-gaussian}

In lecture the target is the law of the terminal log-return $x_T=\log(S_T/S)$, which is
$\mathcal N(\nu T,\sigma^2T)$ under $\Q$, and the sampling law is
$\mathcal N\big((\nu+\alpha)T,\beta^2\sigma^2T\big)$: $\alpha$ shifts the drift of the
log-price and $\beta$ rescales its volatility. For $\beta=1$ the weight is
\begin{equation}\label{eq:is-weight-course}
  L(x_T)=\exp\Big(\frac{\alpha\,[(\alpha+2\nu)T-2x_T]}{2\sigma^2}\Big),
\end{equation}
as one checks by expanding the two Gaussian exponents. In standard units the target is
$\mathcal N(m,s^2)$ with $m=\nu T$, $s=\sigma\sqrt T$, and the sampling law is
$\mathcal N(m+a,\beta^2s^2)$ with $a=\alpha T$, a shift of $c=a/s$ standard deviations.

\begin{proposition}[Mean shift and scale]\label{prop:is-gaussian}
Let the target law be $\mathcal N(m,s^2)$ and the sampling law $\mathcal N(m+a,\beta^2s^2)$,
with $s>0$, $\beta>0$, $a\in\R$, and $c=a/s$.
\begin{enumerate}[label=(\roman*)]
  \item The weight has second moment
    \[
      \Et[L^2]=\EQ[L]=\frac{\beta^2}{\sqrt{2\beta^2-1}}\exp\Big(\frac{c^2}{2\beta^2-1}\Big)
      \quad\text{if }\beta>\tfrac1{\sqrt2},\qquad \Et[L^2]=\infty\quad\text{if }\beta\le\tfrac1{\sqrt2}.
    \]
    For $\beta=1$ this is $e^{c^2}$.
  \item (Mirror formula.) If $\beta>1/\sqrt2$, then for every Borel $f$,
    \begin{equation}\label{eq:is-mirror}
      \Et\big[f(X)^2L(X)^2\big]=\Et[L^2]\;\E\big[f(\tilde X)^2\big],\qquad
      \tilde X\sim\mathcal N\Big(m-\frac{a}{2\beta^2-1},\ \frac{\beta^2s^2}{2\beta^2-1}\Big).
    \end{equation}
  \item The importance-sampling estimator of $\EQ[f(X)]$ has finite variance for every $f$
    of at most exponential growth ($|f(x)|\le Ce^{\kappa|x|}$ for some constants $C,\kappa$)
    if and only if $\beta>1/\sqrt2$.
\end{enumerate}
\end{proposition}

\begin{proof}
Write $u=x-m$ and $\gamma=2\beta^2-1$. The target density is
$(2\pi s^2)^{-1/2}e^{-u^2/(2s^2)}$ and the sampling density
$(2\pi\beta^2s^2)^{-1/2}e^{-(u-a)^2/(2\beta^2s^2)}$, so
\[
  \frac{q(x)^2}{\tilde q(x)}=\frac{\beta}{s\sqrt{2\pi}}\exp\Big(-\frac{u^2}{s^2}+\frac{(u-a)^2}{2\beta^2s^2}\Big)
  =\frac{\beta}{s\sqrt{2\pi}}\exp\Big(-\frac{\gamma u^2+2au-a^2}{2\beta^2s^2}\Big).
\]
If $\gamma>0$, complete the square, $\gamma u^2+2au=\gamma(u+a/\gamma)^2-a^2/\gamma$, and use
$1+1/\gamma=2\beta^2/\gamma$:
\[
  \frac{q(x)^2}{\tilde q(x)}=\frac{\beta}{s\sqrt{2\pi}}\,e^{c^2/\gamma}
  \exp\Big(-\frac{(x-m+a/\gamma)^2}{2\beta^2s^2/\gamma}\Big)
  =\frac{\beta^2}{\sqrt\gamma}\,e^{c^2/\gamma}\;\tilde p(x),
\]
where $\tilde p$ is the density of $\tilde X$. By \eqref{eq:is-variance},
$\Et[f^2L^2]=\int f^2q^2/\tilde q\,\dd x$; taking $f\equiv1$ gives (i) and a general $f$ gives
(ii). If $\gamma<0$ the exponent tends to $+\infty$ quadratically as $|u|\to\infty$, and if
$\gamma=0$ it is affine in $u$, $(a^2-2au)/(2\beta^2s^2)$, whose exponential is not integrable
over $\R$; in both cases $\Et[L^2]=\infty$.
(iii) If $\beta>1/\sqrt2$, then $\E[f(\tilde X)^2]\le C^2\E[e^{2\kappa|\tilde X|}]<\infty$
for a Gaussian $\tilde X$, so (ii) gives a finite variance. If $\beta\le1/\sqrt2$, the payoff
$f\equiv1$, of exponential growth with $\kappa=0$, has infinite variance by (i).
\end{proof}

The mirror formula has a simple reading for $\beta=1$:
$\Et[(fL)^2]=e^{c^2}\,\EQ\big[f(X-a)^2\big]$. Shifting the sampling law by $a$ towards the
payoff costs the factor $e^{c^2}$, and in exchange the second moment is that of the payoff
evaluated at the target shifted by $a$ \emph{the other way}, where a call or a digital is
much smaller. The best shift balances the two. For the call,
$\E[((Se^{Y}-K)^+)^2]$ with $Y\sim\mathcal N(\mu_Y,v)$ equals
\[
  S^2e^{2\mu_Y+2v}N(\delta_2)-2KSe^{\mu_Y+v/2}N(\delta_1)+K^2N(\delta_0),\qquad
  \delta_j=\frac{\mu_Y+jv-\log(K/S)}{\sqrt v},
\]
so the variance of the importance-sampling estimator is an explicit function of $(\alpha,\beta)$,
although its minimiser is not. \Cref{fig:is-contours} reproduces the variance-ratio contours of
the course. With $\beta=1$, the optimal shift and the variance ratio are $\alpha^*=0.206$ and
$9.03$ at $K=90$, $\alpha^*=0.258$ and $10.58$ at $K=100$, and $\alpha^*=0.314$ and $14.29$
at $K=110$: an order of magnitude even at the money, and more the further out of the money
the strike. With $\beta$ free at $K=100$, the ratio keeps increasing as $\beta$ decreases,
to $30.9$ at $\beta=0.75$ and $34.6$ as $\beta\downarrow1/\sqrt2$ (with $\alpha^*\approx0.270$):
the best scale sits on the edge of the region where the weight has finite variance.

\begin{figure}[ht]
  \centering
  \includegraphics[width=0.85\textwidth]{ch09-contours}
  \caption{Variance of plain Monte Carlo over variance of importance sampling for the test-case
    call, computed exactly from \cref{prop:is-gaussian}. (a) Against the strike and the drift
    shift, $\beta=1$, with the optimal shift $\alpha^*(K)$. (b) Against the scale and the shift
    at $K=100$; to the left of $\beta=1/\sqrt2$ the variance is infinite.}
  \label{fig:is-contours}
\end{figure}

\begin{remark}[The share measure and the limits of $\alpha^*$]\label{rem:is-share}
With $\alpha=\sigma^2$ and $\beta=1$, \eqref{eq:is-weight-course} gives
$S_TL=S\exp\big(x_T+(\sigma^2/2+\nu)T-x_T\big)=Se^{(r-y)T}$ for every sample: the forward is
estimated with zero variance. The sampling law, under which $\log S_T$ has drift
$r-y+\sigma^2/2$, is the share measure of \cref{prop:ma-numeraire}, and its density with
respect to $\Q$ is $S_T/\EQ[S_T]$, the zero-variance density of \cref{prop:is-zero-variance}
for $f(x)=Se^x$. Numerically $\alpha^*\to\sigma^2$ as $K\to0$
($\alpha^*=0.0501$ at $K=20$, $0.0404$ at $K=1$), and as $K\to\infty$, heuristically,
$c^*=d+\frac{3}{2d}+o(1/d)$ (checked numerically in our code), the leading terms of the
lecture's expansion.
\end{remark}

\begin{remark}[Thin tails are dangerous, fat tails are safe]\label{rem:is-tails}
For every $\beta<1$ the weight $L(x)$ grows like $\exp(\kappa x^2)$, $\kappa>0$, in both tails,
while $L\tilde q=q$ always decays. What fails for $\beta<1/\sqrt2$ is the integrand of the
second moment, $L^2\tilde q=q^2/\tilde q$, which then grows like $\exp(\kappa'x^2)$. At
$\beta=1/\sqrt2$ exactly, the weight has infinite variance, but the call does not: $q^2/\tilde q$ is then the exponential of an affine function
of $x$, and $\Et[(fL)^2]<\infty$ if and only if $\alpha>\sigma^2$ (compare the growth $e^{2x}$
of $f^2$ with the decay of $q^2/\tilde q$). For $\beta>1$, on the other hand, the exponent $-u^2/(2s^2)+(u-a)^2/(2\beta^2s^2)$ of $L$
(with $u=x-m$) is a concave quadratic in $u$, so the weight is \emph{bounded}, and
$\sigma^2_{\mathrm{is}}+\mu^2=\EQ[f^2L]\le\sup L\,\EQ[f^2]$: whatever is finite under $\Q$
stays finite. (For $\beta=1$ the exponent is affine and $L$ is bounded on a half-line only.)
The same computation with $q^4/\tilde q^3$ shows
that the fourth moment of the call estimator is infinite whenever $\beta^2<3/4$ (and, for the
$n$-th moment of the weight, $\beta^2>1-1/n$; \exref{ex:is-moments}). For
$1/\sqrt2<\beta<\sqrt3/2$, then, $\hat s_n$ converges, but its own fluctuations have infinite
variance, and the error bars are least trustworthy exactly where the course's optimum lies.
\Cref{fig:is-degeneracy}(b) shows the running $\hat s_n$ for $\beta=1$, $0.75$ and $0.55$:
the first two settle at their exact values $3.88$ and $2.27$ (against $12.57$ for plain
Monte Carlo), the third never does.
\end{remark}


\demo{09-importance-sampling.ipynb}{move $\alpha$, $\beta$ and $K$; watch an infinite variance.}

\section{Exponential tilting and large deviations}\label{sec:is-tilting}

The mean shift is a special case of a general recipe. Let $X$ be a real random variable whose
cumulant generating function $\psi(\theta)=\log\EQ[e^{\theta X}]$ is finite on an open
interval $\Theta\ni0$.

\begin{definition}[Exponential tilting]\label{def:is-tilting}
For $\theta\in\Theta$, the \emph{exponentially tilted} (or \emph{Esscher}) law of $X$ has
density $q_\theta(x)=e^{\theta x-\psi(\theta)}q(x)$; its weight is
$L=e^{-\theta X+\psi(\theta)}$. For a sum $S_n=X_1+\dots+X_n$ of independent copies, each
tilted, the weight is $e^{-\theta S_n+n\psi(\theta)}$.
\end{definition}

On $\Theta$, $\psi$ is infinitely differentiable, with $\psi'(\theta)=\E_\theta[X]$ and
$\psi''(\theta)=\Var_\theta(X)$, the mean and variance under the tilted law: for $\theta$ in
$\Theta$ and $\varepsilon>0$ with $\theta\pm\varepsilon\in\Theta$,
$|x|^je^{\theta x}\le C_j(e^{(\theta+\varepsilon)x}+e^{(\theta-\varepsilon)x})$, so dominated
convergence allows differentiation under $\EQ$. If $X$ is not almost surely constant,
$\psi''>0$ and $\psi'$ is strictly increasing: tilting with $\theta>0$ raises the mean.
Examples: for $\mathcal N(m,s^2)$, $\psi(\theta)=m\theta+s^2\theta^2/2$ and the tilted law is
$\mathcal N(m+\theta s^2,s^2)$, the mean shift of \cref{sec:is-gaussian} with $a=\theta s^2$;
for a Poisson count with mean $\lambda$, $\psi(\theta)=\lambda(e^\theta-1)$ and the tilted
law is Poisson with mean $\lambda e^\theta$. Other families, and a law that cannot be tilted,
are in \exref{ex:is-tilts}.

\begin{proposition}[The saddle-point tilt for sums]\label{prop:is-tilting}
Let $X_1,X_2,\dots$ be i.i.d.\ copies of a non-constant $X$ with $\psi$ finite on an open
interval $\Theta\ni0$. Let $a>\EQ[X]$, and suppose that $\psi'(\theta_a)=a$ for some
$\theta_a\in\Theta$. Set $I(a)=a\theta_a-\psi(\theta_a)$ and $p_n=\Q(S_n\ge na)$. Then
$\theta_a>0$ and $I(a)>0$, and:
\begin{enumerate}[label=(\roman*)]
  \item (Chernoff bound.) $p_n\le e^{-nI(a)}$.
  \item The estimator $Y=\1_{\{S_n\ge na\}}e^{-\theta_aS_n+n\psi(\theta_a)}$, with $S_n$
    sampled under the tilted law, satisfies $0\le Y\le e^{-nI(a)}$ surely, and
    $\Et[Y^2]\le e^{-nI(a)}p_n\le e^{-2nI(a)}$.
  \item $\frac1n\log p_n\to-I(a)$, and
    $\log\Et[Y^2]/\log(p_n^2)\to1$, while plain Monte Carlo has relative variance
    $(1-p_n)/p_n=e^{nI(a)+o(n)}$.
\end{enumerate}
\end{proposition}

\begin{proof}
$\psi'$ is strictly increasing and $\psi'(\theta_a)=a>\EQ[X]=\psi'(0)$, so $\theta_a>0$. By
strict convexity, $0=\psi(0)>\psi(\theta_a)-\theta_a\psi'(\theta_a)=-I(a)$.
(ii) On $\{S_n\ge na\}$, $\theta_a>0$ gives $e^{-\theta_aS_n}\le e^{-n\theta_aa}$, so
$Y\le e^{-n(a\theta_a-\psi(\theta_a))}=e^{-nI(a)}$; hence $Y^2\le e^{-nI(a)}Y$ and
$\Et[Y^2]\le e^{-nI(a)}\Et[Y]=e^{-nI(a)}p_n$, by \cref{prop:is-likelihood-ratio}(i).
(i) $p_n=\Et[Y]\le e^{-nI(a)}$.
(iii) The function $\theta\mapsto a\theta-\psi(\theta)$ is concave with a critical point at
$\theta_a$, so $I(a)=\sup_\theta(a\theta-\psi(\theta))$, the Legendre transform of $\psi$ at
$a$. By Cramér's theorem \citep[Theorem~2.2.3]{dembo-zeitouni},
$\liminf_n\frac1n\log\Q(S_n/n>a)\ge-\inf_{x>a}I(x)$. The function $I$ is convex, $I\ge0$ (take $\theta=0$) and
$I(\EQ X)=0$ (Jensen: $\psi(\theta)\ge\theta\,\EQ X$), so $I$ is nondecreasing on
$[\EQ X,\infty)$; it is finite near $a$ (every $x=\psi'(\theta)$ with $\theta$ near $\theta_a$
is attained), hence continuous at $a$. So $\inf_{x>a}I(x)=I(a)$ and the right side is $-I(a)$; with (i), $\frac1n\log p_n\to-I(a)$. Jensen's inequality
gives $\Et[Y^2]\ge p_n^2$, and (ii) gives $\Et[Y^2]\le e^{-nI(a)}p_n$; dividing the logarithms
by $\log p_n^2<0$,
\[
  1\ \ge\ \frac{\log\Et[Y^2]}{\log p_n^2}\ \ge\ \frac{-nI(a)+\log p_n}{2\log p_n}\ \longrightarrow\ 1 .\qedhere
\]
\end{proof}

The tilt $\theta_a$ is a \emph{saddle point}: it makes the rare level the mean of the
sampling law, $\E_{\theta_a}[S_n/n]=a$, and no more. Given the rare event, the summands behave approximately like tilted
variables (\citealp{asmussen-glynn}, Sections~VI.2 and~VI.5). As a jump-count example, let $S_n$ be the
number of jumps of a Poisson process with intensity $2$ per year over $n=10$ years, and
$p_n=\Q(S_n\ge40)=5.32\times10^{-5}$. Then $\theta_a=\log2$, $I(a)=4\log2-2=0.773$, the bound
is $e^{-nI(a)}=4.41\times10^{-4}$, and plain Monte Carlo has relative error $137$ per sample.
Tilted counts (Poisson with mean $4$ per year) give $2.17$ (our implementation,
\texttt{default\_rng}, seed $7$, $10^5$ samples): $(137/2.17)^2\approx4000$ times fewer samples.

\section{Path space: Girsanov's theorem and the optimal drift}\label{sec:is-path}

A simulated path is a function of a vector $Z\in\R^m$ of independent normals, one per date
(or one per date and factor). Shifting $Z$ to $\mathcal N(\mu,I_m)$ multiplies $m$
independent one-dimensional weights, so by \cref{prop:is-gaussian} with $\beta=1$
\begin{equation}\label{eq:is-path-weight}
  L(Z)=\exp\big(-\mu^\top Z+\tfrac12|\mu|^2\big),\qquad \Et[L^2]=e^{|\mu|^2}.
\end{equation}
The continuous-time meaning of this weight is Girsanov's theorem.

\begin{theorem}[Path-space importance sampling]\label{thm:is-girsanov}
Let $B$ be a standard Brownian motion on $[0,T]$ under $\Qt$, with a filtration for it, and let
$\theta$ be an adapted process that is either bounded or deterministic with
$\int_0^T\theta_t^2\,\dd t<\infty$. Set
\[
  \tilde W_t=B_t+\int_0^t\theta_u\,\dd u,\qquad
  L_T=\exp\Big(-\int_0^T\theta_u\,\dd B_u-\frac12\int_0^T\theta_u^2\,\dd u\Big).
\]
\begin{enumerate}[label=(\roman*)]
  \item For every Borel functional $G$ on $C[0,T]$ with $\EQ|G(W)|<\infty$ for a
    $\Q$-Brownian motion $W$: $\ \EQ[G(W)]=\Et\big[G(\tilde W)\,L_T\big]$.
  \item If $\theta$ is deterministic, $\Et[L_T^2]=\exp\big(\int_0^T\theta_u^2\,\dd u\big)$.
  \item If $\theta_u=\mu_j/\sqrt{\Delta t}$ on $(t_{j-1},t_j]$, $t_j=j\Delta t$, $\Delta t=T/m$,
    then $L_T=\exp(-\mu^\top Z+\frac12|\mu|^2)$ with
    $Z_j=(\tilde W_{t_j}-\tilde W_{t_{j-1}})/\sqrt{\Delta t}$, and the $Z_j$ are independent
    $\mathcal N(\mu_j,1)$ under $\Qt$: the discrete mean shift \eqref{eq:is-path-weight} is
    this change of measure exactly, with no discretisation error in the weight.
  \item Under $\Qt$, $S_t=S\exp(\nu t+\sigma\tilde W_t)$ solves
    $\dd S_t=(r-y+\sigma\theta_t)S_t\,\dd t+\sigma S_t\,\dd B_t$.
\end{enumerate}
\end{theorem}

A feedback drift $\theta_t=h(t,\tilde W_t)$ makes $\tilde W$ the solution of
$\dd\tilde W=h(t,\tilde W)\,\dd t+\dd B$, and $\theta$ is adapted only if this equation has a
strong solution, for instance when $h$ is Lipschitz in $x$; if $h$ is also bounded, the theorem
applies.

\begin{proof}
(i) We use Girsanov's theorem in the form of \textcite{shreve2}, Theorem~5.2.3, with
$\Theta=\theta$ and $W=B$: if $\Et\int_0^T\theta_u^2L_u^2\,\dd u<\infty$, where $L_u$ is the
same exponential on $[0,u]$, then $\Et[L_T]=1$ and $\tilde W$ is a Brownian motion under the
probability $\hat\Q$ with $\dd\hat\Q=L_T\,\dd\Qt$. The hypothesis holds: if $|\theta|\le C$,
then $L_u^2=M_u\exp(\int_0^u\theta^2)\le M_ue^{C^2u}$ with
$M_u=\exp(-2\int_0^u\theta\,\dd B-2\int_0^u\theta^2)$ a positive local martingale, hence a
supermartingale with $\Et M_u\le1$; if $\theta$ is deterministic, (ii) below on $[0,u]$ gives
$\Et L_u^2=e^{\int_0^u\theta^2}$. The law of a Brownian motion on $C[0,T]$ is determined by
its finite-dimensional distributions, so $G(\tilde W)$ has the same law under $\hat\Q$ as
$G(W)$ under $\Q$, and $\EQ[G(W)]=\E_{\hat\Q}[G(\tilde W)]=\Et[L_TG(\tilde W)]$.
(ii) For deterministic $\theta$, $\int_0^T\theta\,\dd B\sim\mathcal N(0,v)$ with
$v=\int_0^T\theta^2$ \citep[Theorem~4.4.9]{shreve2}, so
$\Et[L_T^2]=\Et[e^{-2\int\theta\,\dd B}]e^{-v}=e^{2v}e^{-v}$.
(iii) $\int_0^T\theta\,\dd B=\sum_j\mu_j(B_{t_j}-B_{t_{j-1}})/\sqrt{\Delta t}=\sum_j\mu_j(Z_j-\mu_j)$,
because $\tilde W_{t_j}-\tilde W_{t_{j-1}}=B_{t_j}-B_{t_{j-1}}+\mu_j\sqrt{\Delta t}$, and
$\frac12\int\theta^2=\frac12|\mu|^2$; so $L_T=\exp(-\mu^\top Z+|\mu|^2-\frac12|\mu|^2)$. The
Brownian increments are independent $\mathcal N(0,\Delta t)$ under $\Qt$.
(iv) It\^o's formula: $\dd S=S\big((\nu+\frac{\sigma^2}2)\,\dd t+\sigma\,\dd\tilde W\big)$ and
$\dd\tilde W=\dd B+\theta\,\dd t$.
\end{proof}

So a drift added to the Brownian motion is a drift $\sigma\theta$ added to the asset: to price
an out-of-the-money call, simulate a market that rallies, and weight. A volatility cannot be
changed this way. If each of $m$ increments is rescaled by $\beta\ne1$, the weight has
$\Et[L^2]=\big(\beta^2/\sqrt{2\beta^2-1}\big)^m$ by \cref{prop:is-gaussian}, and
$\beta^4>2\beta^2-1$ makes this explode as the grid is refined; in the limit the two path
laws are mutually singular, because the realised quadratic variation of $\log S$ is
$\sigma^2T$ almost surely under one and $\beta^2\sigma^2T$ under the other. The course's
$\beta$ is a device for one coordinate, such as the terminal value.

\begin{remark}[The zero-variance law in path space is a feedback drift]\label{rem:is-feedback}
Let $\dd X_t=b(t,X_t)\,\dd t+\sigma(t,X_t)\,\dd W_t$ under $\Q$, with $b$ and $\sigma$
continuous, let $g>0$, and suppose that $u\in C^{1,2}([0,T)\times\R)$ is positive, continuous
on $[0,T]\times\R$, solves $u_t+bu_x+\frac12\sigma^2u_{xx}=0$ with $u(T,\cdot)=g$, and
that $\int\sigma u_x(t,X_t)\,\dd W_t$ is a martingale. It\^o's formula and the equation give
$\dd u(t,X_t)=\sigma u_x\,\dd W_t$, so $u(t,X_t)=\EQ[g(X_T)\mid\F_t]$ is the price process (the
Feynman--Kac representation of \cref{thm:sm-feynman-kac}, with zero rate), and It\^o's formula
for $\log u$ gives
\[
  g(X_T)=u(0,X_0)\exp\Big(\int_0^T\theta^*_t\,\dd W_t-\frac12\int_0^T(\theta^*_t)^2\,\dd t\Big),\qquad
  \theta^*_t=\sigma(t,X_t)\,\partial_x\log u(t,X_t).
\]
The exponential is the martingale $u(t,X_t)/u(0,X_0)$, so Girsanov's theorem for a martingale
density \citep[Theorem~3.5.1]{karatzas-shreve} applies: the zero-variance law of
\cref{prop:is-zero-variance}, with density $g(X_T)/u(0,X_0)$ with respect to $\Q$, is the
change of measure under which $W$ acquires the drift $\theta^*$. It needs $u$, the price
function itself; the drift of the next proposition freezes it along a single path.
\end{remark}

\begin{proposition}[The Glasserman--Heidelberger--Shahabuddin drift]\label{prop:is-optimal-drift}
Let $G:\R^m\to[0,\infty)$ be Borel with $p=\E[G(Z)]\in(0,\infty)$ for $Z\sim\mathcal N(0,I_m)$,
let $F=\log G$ (with $F=-\infty$ where $G=0$), and $H(z)=F(z)-\frac12|z|^2$. For $\mu\in\R^m$
let $Y_\mu=G(Z)\exp(-\mu^\top Z+\frac12|\mu|^2)$ with $Z\sim\mathcal N(\mu,I_m)$ under $\Qt$.
\begin{enumerate}[label=(\roman*)]
  \item $\Et[Y_\mu]=p$ and $\Et[Y_\mu^2]=\E\big[G(Z)^2e^{-\mu^\top Z+|\mu|^2/2}\big]$ for
    $Z\sim\mathcal N(0,I_m)$.
  \item If $G(z)=e^{\kappa+v^\top z}$, then $v$ is the unique maximiser of $H$ and
    $Y_v=p$ surely: zero variance.
  \item If $F$ is concave and $\mu$ maximises $H$, then $F(z)\le F(\mu)+\mu^\top(z-\mu)$ for
    all $z$. Consequently $0\le Y_\mu\le e^{H(\mu)}$ surely, $p\le e^{H(\mu)}$, and
    $\Et[Y_\mu^2]/p^2\le e^{H(\mu)}/p$.
\end{enumerate}
At a maximiser where $G$ is positive and differentiable, $\nabla F(\mu)=\mu$, that is
$\nabla G(\mu)=G(\mu)\,\mu$.
\end{proposition}

\begin{proof}
(i) is \cref{prop:is-likelihood-ratio} with \eqref{eq:is-path-weight}.
(ii) $H(z)=\kappa+\frac12|v|^2-\frac12|z-v|^2$ is uniquely maximised at $v$, and
$Y_v=e^{\kappa+v^\top Z-v^\top Z+|v|^2/2}=e^{\kappa+|v|^2/2}=\E[e^{\kappa+v^\top Z}]=p$.
(iii) $H(\mu)$ is finite since $p>0$ forces $G>0$ somewhere. Fix $z$ with $F(z)>-\infty$
(otherwise there is nothing to prove) and $t\in(0,1]$. Concavity gives
$F(\mu+t(z-\mu))\ge(1-t)F(\mu)+tF(z)$, and maximality gives
$H(\mu+t(z-\mu))\le H(\mu)$. Together,
\[
  t\big(F(z)-F(\mu)\big)\le\tfrac12|\mu+t(z-\mu)|^2-\tfrac12|\mu|^2=t\,\mu^\top(z-\mu)+\tfrac{t^2}2|z-\mu|^2 .
\]
Divide by $t$ and let $t\downarrow0$. Then
$Y_\mu=\exp\big(F(Z)-\mu^\top Z+\frac12|\mu|^2\big)\le\exp\big(F(\mu)-\frac12|\mu|^2\big)=e^{H(\mu)}$,
so $p=\Et[Y_\mu]\le e^{H(\mu)}$ and $\Et[Y_\mu^2]\le e^{H(\mu)}\Et[Y_\mu]=e^{H(\mu)}p$.
The first-order condition is $\nabla H(\mu)=0$.
\end{proof}

The drift $\mu$ is the \emph{most important path}: the point that maximises payoff times
density, $G(z)e^{-|z|^2/2}$. Writing
$Y_\mu=e^{H(\mu)}\exp\big(F(Z)-F(\mu)-\nabla F(\mu)^\top(Z-\mu)\big)$ shows what the shift
achieves: the linear part of $\log G$ around $\mu$ is removed exactly, and only its curvature
contributes variance. The drift is asymptotically optimal in a small-noise limit
(\cref{rem:is-small-noise}); \textcite{glasserman-heidelberger-shahabuddin} combine it with
stratification along the directions where the curvature is largest (\cref{ch:10-stratification}). A call on a lognormal asset has a concave $\log G$,
since $u\mapsto\log(e^u-K)$ is concave on $\{e^u>K\}$; so do a digital on a convex set and the
geometric Asian call (\exref{ex:is-geometric}). For the arithmetic Asian, $\log G$ need not be
concave, and the drift is found numerically.

\begin{remark}[Small-noise optimality]\label{rem:is-small-noise}
For $\varepsilon>0$ let $p_\varepsilon=\E\big[\exp\big(F(\sqrt\varepsilon Z)/\varepsilon\big)\big]$,
$Z\sim\mathcal N(0,I_m)$: noise of size $\sqrt\varepsilon$ and a log-payoff of order
$1/\varepsilon$, the large-deviations regime. Sample with the drift $\mu/\sqrt\varepsilon$, where
$\mu$ maximises $H$, and call the estimator $Y_\varepsilon$. \emph{Assume} that $F$ is concave
and finite and continuous at $\mu$. Then $\varepsilon\log p_\varepsilon\to H(\mu)$ and
$\varepsilon\log\Et[Y_\varepsilon^2]\to2H(\mu)$: the estimator is logarithmically efficient in
$\varepsilon$. Indeed $F_\varepsilon(z)=F(\sqrt\varepsilon z)/\varepsilon$ is concave with
$F_\varepsilon(z)-\frac12|z|^2=H(\sqrt\varepsilon z)/\varepsilon$, maximised at $\mu/\sqrt\varepsilon$, so
(iii) gives $p_\varepsilon\le e^{H(\mu)/\varepsilon}$ and $\Et[Y_\varepsilon^2]\le e^{H(\mu)/\varepsilon}p_\varepsilon$.
Conversely, given $\delta>0$ take $r\in(0,\delta]$ with $F\ge F(\mu)-\delta$ on the ball
$B(\mu,r)$; shifting $Z$ by $\mu/\sqrt\varepsilon$ as in \eqref{eq:is-path-weight},
$\Q(\sqrt\varepsilon Z\in B(\mu,r))\ge e^{-(|\mu|^2/2+|\mu|r)/\varepsilon}\,\Q(|Z|<r/\sqrt\varepsilon)$, so
$\liminf\varepsilon\log p_\varepsilon\ge H(\mu)-\delta(1+|\mu|)$. Jensen's inequality
$\Et[Y_\varepsilon^2]\ge p_\varepsilon^2$ closes the argument. \textcite{glasserman-heidelberger-shahabuddin}
establish this kind of optimality for smooth payoffs that need not be log-concave; see
\textcite{glasserman-mc}, Section~4.6.2, for the statement and its hypotheses.
\end{remark}

\begin{example}[Arithmetic Asian call]\label{exa:is-asian}
Take the test case with $m=16$ equally spaced fixings $t_j=j/16$ and payoff
$e^{-rT}\big(\frac1{16}\sum_jS_{t_j}-K\big)^+$. Maximising $H$ by a quasi-Newton method from a
point where the payoff is positive gives the drifts of \cref{fig:is-asian}(a): they decrease
along the path, because an early normal moves every later fixing of the average, and their
size $|\mu|$ is $1.10$, $2.11$, $3.17$ and $4.13$ for $K=100$, $120$, $140$, $160$. The
variance falls by factors $9.2$, $42$, $618$ and $1.9\times10^4$ (\cref{tab:is-lab}).
\end{example}

\begin{figure}[ht]
  \centering
  \includegraphics[width=0.85\textwidth]{ch09-asian}
  \caption{Arithmetic Asian call, test case, $16$ fixings. (a) The GHS drift
    (lines) and the cross-entropy drift after $8$ iterations of $10^4$ paths (dots;
    \cref{sec:is-ce}). (b) Variance reduction factors of the control variate, of importance
    sampling and of both together (\cref{tab:is-lab}).}
  \label{fig:is-asian}
\end{figure}

\section{Rare events and relative error}\label{sec:is-rare}

When the price $p$ is small, the absolute error says little; what matters is the
\emph{relative error per sample} $\sqrt{\Var(Y)}/p=\sqrt{\Et[Y^2]/p^2-1}$. A relative
precision $\varepsilon$ at $95\%$ confidence needs about $(1.96/\varepsilon)^2\,\Et[Y^2]/p^2$
samples. For plain Monte Carlo of a probability, $\Et[Y^2]/p^2=1/p$.

\begin{definition}[Efficiency for rare events]\label{def:is-efficiency}
Let $Y_d$ be unbiased estimators of $p_d$, with $p_d\to0$ as $d\to\infty$. The family has
\emph{bounded relative error} if $\limsup_{d\to\infty}\Et[Y_d^2]/p_d^2<\infty$, and is
\emph{logarithmically efficient} if $\log\Et[Y_d^2]/\log(p_d^2)\to1$
\citep[Section~VI.1]{asmussen-glynn}.
\end{definition}

Since $\Et[Y_d^2]\ge p_d^2$, the ratio of logarithms never exceeds $1$ once $p_d<1$:
logarithmic efficiency says that the second moment decays at the best possible exponential
rate, while the relative error may still grow, but more slowly than any power of $1/p_d$.
\Cref{prop:is-tilting} is an example.

\begin{lemma}[Mills ratio]\label{lem:is-mills}
For $x>0$, $\ \dfrac{x}{1+x^2}N'(x)\le N(-x)\le\dfrac{N'(x)}{x}$.
\end{lemma}

\begin{proof}
$N(-x)=\int_x^\infty N'(t)\,\dd t\le\int_x^\infty\frac tx N'(t)\,\dd t=\frac{N'(x)}x$, since
$(N')'(t)=-tN'(t)$. And $\frac{\dd}{\dd t}\big(-N'(t)/t\big)=(1+t^{-2})N'(t)$, so
$N'(x)/x=\int_x^\infty(1+t^{-2})N'(t)\,\dd t\le(1+x^{-2})N(-x)$.
\end{proof}

\begin{proposition}[Deep out-of-the-money digitals]\label{prop:is-rare}
Let $d>0$ and $p_d=\Q(Z>d)=N(-d)$: the undiscounted value of a digital call with
$d=-d_-$.
\begin{enumerate}[label=(\roman*)]
  \item Plain Monte Carlo has $\Et[Y^2]/p_d^2=1/p_d$, between $\sqrt{2\pi}\,de^{d^2/2}$ and
    $\sqrt{2\pi}\,\frac{1+d^2}{d}e^{d^2/2}$: exponential growth in $d^2$.
  \item The mean shift to the strike, $Z\sim\mathcal N(d,1)$ under $\Qt$, has
    \[
      \frac{2\sqrt{2\pi}\,d^3}{1+4d^2}\ \le\ \frac{\Et[Y^2]}{p_d^2}=\frac{e^{d^2}N(-2d)}{N(-d)^2}\ \le\ \frac{\sqrt{2\pi}\,(1+d^2)^2}{2d^3},
    \]
    both bounds being $\sqrt{\pi/2}\,d\,(1+o(1))$: logarithmically efficient, but the relative
    error grows like $(\pi/2)^{1/4}\sqrt d$.
  \item Sampling $Z=d+E/d$ with $E$ standard exponential gives $Y\le N'(d)/d$ surely and
    $\Et[Y^2]/p_d^2\le(1+d^2)/d^2$: the relative error per sample is at most $1/d$, so it
    vanishes as $d\to\infty$.
\end{enumerate}
\end{proposition}

\begin{proof}
(i) is \cref{lem:is-mills} applied to $1/p_d$. (ii) By the mirror formula with $\beta=1$ and
$a=c=d$, $\Et[Y^2]=e^{d^2}\Q(Z-d>d)=e^{d^2}N(-2d)$. \Cref{lem:is-mills} at $2d$ and at $d$,
with $e^{d^2}N'(2d)=e^{-d^2}/\sqrt{2\pi}$ and $N'(d)^2=e^{-d^2}/(2\pi)$, gives
\[
  \frac{e^{d^2}N(-2d)}{N(-d)^2}\le\frac{e^{-d^2}/(2d\sqrt{2\pi})}{d^2e^{-d^2}/(2\pi(1+d^2)^2)},\qquad
  \frac{e^{d^2}N(-2d)}{N(-d)^2}\ge\frac{2d\,e^{-d^2}/((1+4d^2)\sqrt{2\pi})}{e^{-d^2}/(2\pi d^2)},
\]
which simplify to the stated bounds. Logarithmic efficiency follows because
$\log p_d^2=-d^2+O(\log d)$ and $\log\Et[Y^2]=-d^2+O(\log d)$.
(iii) The sampling density is $\tilde q(z)=de^{-d(z-d)}$ on $(d,\infty)$, which contains the
support of $f=\1_{(d,\infty)}$. Since $N'(z)=N'(d)e^{-d(z-d)-(z-d)^2/2}$, the weight is
$L(z)=N'(d)e^{-(z-d)^2/2}/d\le N'(d)/d$, and $Y=L(Z)$. Then
$\Et[Y^2]\le\frac{N'(d)}d\Et[Y]=\frac{N'(d)}dp_d$, and \cref{lem:is-mills} gives
$N'(d)/(dp_d)\le(1+d^2)/d^2$.
\end{proof}

The exponential law in (iii) is, to first order, the conditional law of $Z-d$ given $Z>d$:
it approximates the zero-variance law of \cref{prop:is-zero-variance} in shape, not just in
location, which a Gaussian cannot do. For the call, with the same $d$, the relative second
moments grow like $2\sqrt{2\pi}\,de^{d^2/2}$ for plain Monte Carlo and like $\sqrt{\pi/8}\,d$
for the shift to the strike (\exref{ex:is-call-asymptotics}). For the digital, written as $G=\1_{[d,\infty)}$ (same price), $\log G$ is concave, the GHS drift
is the boundary point $d$, and the bound of \cref{prop:is-optimal-drift}(iii) reads
$\Et[Y^2]/p_d^2\le e^{-d^2/2}/N(-d)\le\sqrt{2\pi}(1+d^2)/d$, about twice the exact value.

\Cref{tab:is-drifts} compares, at two strikes, the shifts chosen by three criteria: the GHS
drift, the variance-optimal shift, and the cross-entropy shift of \cref{sec:is-ce}, which for a
digital is $\E[Z\mid Z>d]=N'(d)/N(-d)$. They differ by amounts that shrink as $d$ grows, and
their relative errors are within $3\%$ of each other; the payoff shape matters more than the
criterion. \Cref{fig:is-rare} shows the whole range of strikes.

\begin{table}[ht]
  \centering
  \small
  \begin{tabular}{@{}ll@{\quad}cc@{\quad}cc@{\quad}cc@{\quad}c@{}}
    \toprule
    & & \multicolumn{2}{c}{GHS} & \multicolumn{2}{c}{optimal} & \multicolumn{2}{c}{cross-entropy} & plain \\
    & & shift & rel.\ err. & shift & rel.\ err. & shift & rel.\ err. & rel.\ err. \\
    \midrule
    $K=150$, $d=2.177$ & digital & $2.177$ & $1.587$ & $2.379$ & $1.563$ & $2.531$ & $1.577$ & $8.18$ \\
     & call & $2.581$ & $1.010$ & $2.731$ & $0.998$ & $2.844$ & $1.005$ & $11.43$ \\
    $K=200$, $d=3.616$ & digital & $3.616$ & $2.018$ & $3.747$ & $2.007$ & $3.861$ & $2.015$ & $81.71$ \\
     & call & $3.880$ & $1.312$ & $3.995$ & $1.305$ & $4.092$ & $1.310$ & $115.26$ \\
    \bottomrule
  \end{tabular}
  \caption{Mean shifts (in standard deviations of $Z$) and exact relative errors per sample
    for the test case with strikes $150$ and $200$ (digital prices $0.0146$ and
    $1.48\times10^{-4}$, call prices $0.165$ and $1.52\times10^{-3}$). The exponential sampler
    of \cref{prop:is-rare}(iii) has relative errors $0.234$ and $0.119$ for the digitals.}
  \label{tab:is-drifts}
\end{table}

\begin{figure}[ht]
  \centering
  \includegraphics[width=0.85\textwidth]{ch09-rare}
  \caption{Exact relative error per sample against the strike for the test case: (a) digital
    call, (b) call. The GHS, variance-optimal and cross-entropy shifts are nearly
    indistinguishable; dots are Monte Carlo estimates for the GHS shift ($10^5$ samples).}
  \label{fig:is-rare}
\end{figure}

\section{Adaptive importance sampling: cross-entropy}\label{sec:is-ce}

The variance-optimal parameter minimises $\Et[(fL)^2]$, itself an expectation to be estimated
with weights that may be wild. The \emph{cross-entropy method}
\citep{rubinstein-ce,deboer-ce} replaces the variance by a distance to the zero-variance law
$g^*=fq/\mu$ that, for exponential families, is minimised in closed form.

\begin{proposition}[Cross-entropy for Gaussian mean shifts]\label{prop:is-cross-entropy}
Let $f\ge0$ be Borel, $Z\sim\mathcal N(0,I_m)$ under $\Q$, $\mu=\EQ[f(Z)]\in(0,\infty)$,
$\EQ[f(Z)|Z|^2]<\infty$, and $g^*=fq/\mu$. For $v\in\R^m$ let $\tilde q_v$ be the density of
$\mathcal N(v,I_m)$ and $L_v=q/\tilde q_v$.
\begin{enumerate}[label=(\roman*)]
  \item The cross-entropy $-\int g^*\log\tilde q_v$ equals
    $\frac m2\log(2\pi)+\frac12\E_{g^*}|Z-v^*|^2+\frac12|v-v^*|^2$, with
    $v^*=\E_{g^*}[Z]=\EQ[f(Z)Z]/\mu$. It is uniquely minimised at $v=v^*$; when the
    Kullback--Leibler divergence $\int g^*\log(g^*/\tilde q_v)$ is finite, it differs from the
    cross-entropy by a constant, and $v^*$ minimises it too.
  \item For every $u\in\R^m$, $v^*=\Et_u[fL_uZ]/\Et_u[fL_u]$, and if $Z_1,Z_2,\dots$ are
    independent $\mathcal N(u,I_m)$, the ratio
    $\hat v_n=\sum_if(Z_i)L_u(Z_i)Z_i\big/\sum_if(Z_i)L_u(Z_i)$ converges to $v^*$ almost surely.
\end{enumerate}
\end{proposition}

\begin{proof}
(i) $-\log\tilde q_v(z)=\frac m2\log(2\pi)+\frac12|z-v|^2$, and
$\E_{g^*}|Z-v|^2=\E_{g^*}|Z-v^*|^2+|v-v^*|^2$ because $\E_{g^*}[Z-v^*]=0$; all terms are
finite by the moment assumption. (ii) By \cref{prop:is-likelihood-ratio}(i), applied
coordinatewise to $fZ_j$ and to $f$, $\EQ[fZ]=\Et_u[fL_uZ]$ and $\mu=\Et_u[fL_u]$. The
numerator and denominator of $\hat v_n$, divided by $n$, converge almost surely by the strong
law, and the limit of the denominator is $\mu>0$.
\end{proof}

When $f$ is rare, most $f(Z_i)$ vanish at $u=0$ and $\hat v_n$ is undefined or noisy. The
multilevel algorithm (\citealp{deboer-ce}, Section~2.1 and Algorithm~2.1) climbs to the target
through intermediate levels: sample from $\mathcal N(v_t,I)$, set the level $\gamma_t$ to the
smaller of the target and the $(1-\rho)$-quantile of the simulated scores (for example the
average $A$ of an Asian), update $v_{t+1}=\hat v_n$ with the payoff at level $\gamma_t$, and
stop once $\gamma_t$ has reached the target and $v_t$ has settled. For the Asian example, with
$\rho=0.1$ and $10^4$ paths per iteration, the target level is reached at the second or third
iteration, and after $8$ iterations the ratio of the norms of the cross-entropy and GHS drifts
is $1.23$, $1.13$, $1.07$, $1.05$ for $K=100,120,140,160$ (\cref{fig:is-asian}(a); our
implementation, seed $K$): close in the tail, larger near the money. For the digital, $v^*=\E[Z\mid Z>d]=N'(d)/N(-d)$ overshoots the GHS
drift $d$ by at most $1/d$, by \cref{lem:is-mills}.

Adaptive schemes have risks that a fixed drift does not.
\begin{enumerate}[label=(\alph*)]
  \item \emph{Noisy updates.} Each update is a weighted average. If the current law is far from
    $g^*$, a few samples carry all the weight and successive iterates wander; the effective
    sample size below, computed with the update weights $f(Z_i)L(Z_i)$, is the alarm.
  \item \emph{Reuse bias.} Iteration $t$ samples from $\mathcal N(v_t,I)$ with $v_t$ fixed by
    earlier iterations, so its weighted mean is conditionally unbiased, and an equal-weight pool
    of the iterations' estimates is unbiased (a sum of martingale differences). Bias appears
    when a parameter chosen from a sample is applied to that same sample, for instance
    $f(Z_i+\hat v)e^{-\hat v^\top Z_i-|\hat v|^2/2}$ with $\hat v$ computed from the $Z_i$, or
    when the best of several candidates is reported with its in-sample error: the analogue of
    \cref{prop:cv-estimated-beta}.
  \item \emph{Scale.} The full Gaussian family also updates the covariance, to
    $\mathrm{Cov}_{g^*}(Z)$. A smaller covariance means thinner tails, and
    \cref{prop:is-gaussian} shows how that can make the variance infinite without any sign in
    the sample; in path space, \cref{sec:is-path} shows that rescaling many coordinates
    degenerates. Estimating $m(m+1)/2$ parameters from weighted samples is itself noisy.
\end{enumerate}

\begin{proposition}[Effective sample size]\label{prop:is-ess}
Let $L_1,L_2,\dots\ge0$ be i.i.d.\ weights with $\Et[L]=1$, and
$\mathrm{ESS}_n=\big(\sum_{i\le n}L_i\big)^2/\sum_{i\le n}L_i^2$ (on the event that some
$L_i>0$).
\begin{enumerate}[label=(\roman*)]
  \item $1\le\mathrm{ESS}_n\le n$, with $\mathrm{ESS}_n=n$ if and only if all the $L_i$ are equal.
  \item $\mathrm{ESS}_n/n\to1/\Et[L^2]$ almost surely, with $1/\infty=0$.
  \item For the mean shift $\mathcal N(\mu,I_m)$, $1/\Et[L^2]=e^{-|\mu|^2}$; for a rescaling
    $\mathcal N(0,\beta^2I_m)$ with $\beta>1/\sqrt2$, it is $\big(\sqrt{2\beta^2-1}/\beta^2\big)^m$.
\end{enumerate}
\end{proposition}

\begin{proof}
(i) Cauchy--Schwarz gives $(\sum L_i)^2\le n\sum L_i^2$, with equality iff the $L_i$ are equal;
and $(\sum L_i)^2\ge\sum L_i^2$ for nonnegative terms. (ii) By the strong law,
$\frac1n\sum L_i\to1$; if $\Et[L^2]<\infty$, also $\frac1n\sum L_i^2\to\Et[L^2]$. If
$\Et[L^2]=\infty$, then for every $M$, $\frac1n\sum L_i^2\ge\frac1n\sum\min(L_i^2,M)\to\Et[\min(L^2,M)]$,
which tends to $\infty$ as $M\to\infty$ by monotone convergence. (iii) The weights are products
of $m$ independent one-dimensional weights; use \cref{prop:is-gaussian}(i).
\end{proof}

The effective sample size of a mean shift depends on $|\mu|^2$ only. A drift of fixed size
spread over many dates, like a Girsanov drift with fixed $\int\theta^2$ on a finer and finer
grid, does not degenerate as the dimension grows; a fixed shift $c$ per coordinate gives
$ne^{-mc^2}$, and a rescaling decays geometrically in $m$ (\cref{fig:is-degeneracy}(a)). The
effective sample size describes the weights, not the estimator of a given payoff. For the
Asian at $K=160$, $e^{-|\mu|^2}\approx4\times10^{-8}$, and yet the estimator is excellent: the
weights are deliberately small on the paths that pay and large on paths that do not. For one
payoff the diagnostic is the relative standard error of $fL$; the effective sample size matters
when one weighted sample prices many payoffs, and inside the adaptive updates.

\begin{figure}[ht]
  \centering
  \includegraphics[width=0.85\textwidth]{ch09-degeneracy}
  \caption{(a) $\mathrm{ESS}/n$ against the dimension (\cref{prop:is-ess}): lines exact, dots
    from $10^4$ draws, which overestimate it when the weights are heavy-tailed. (b) Running
    sample standard deviation of the call estimator with $\alpha=0.27$: for $\beta=1$ and $0.75$
    it settles at the exact values (dashed); for $\beta=0.55<1/\sqrt2$ the variance is infinite
    and three seeds never settle (our implementation, \texttt{default\_rng}, seeds $1$--$3$).}
  \label{fig:is-degeneracy}
\end{figure}

\section{Importance sampling with control variates}\label{sec:is-cv}

Importance sampling decides where to sample; a control variate removes what is predictable
in the samples. A control keeps its known mean after weighting.

\begin{proposition}[Controls under the sampling law]\label{prop:is-control}
Let $g$ be Borel with known $\gamma=\EQ[g(X)]$, and let $\tilde q$ satisfy
\eqref{eq:is-support} for both $f$ and $g$, with $f(X)L(X)$ and $g(X)L(X)$ in $L^2(\Qt)$.
\begin{enumerate}[label=(\roman*)]
  \item For every $b$, $Y_b=f(X)L(X)-b\big(g(X)L(X)-\gamma\big)$ has $\Et[Y_b]=\mu$. Its
    variance is minimised at $b^*=\widetilde{\Cov}(fL,gL)/\widetilde\Var(gL)$, where it equals
    $(1-\tilde\rho^2)\widetilde\Var(fL)$, $\tilde\rho$ being the correlation of $fL$ and $gL$
    under $\Qt$ (provided $\widetilde\Var(gL)>0$).
  \item With $g\equiv1$ the control is the weight itself, with mean $1$ when $\{q>0\}\subseteq\{\tilde q>0\}$.
    The \emph{self-normalised} estimator $\hat\mu^{\mathrm{sn}}_n=\sum_if(X_i)L(X_i)/\sum_iL(X_i)$
    equals the control-variate estimator with control $L$ and coefficient
    $b=\hat\mu^{\mathrm{sn}}_n$.
\end{enumerate}
\end{proposition}

\begin{proof}
(i) $\Et[gL]=\gamma$ by \cref{prop:is-likelihood-ratio}(i), and the rest is
\cref{thm:cv-optimal} under $\Qt$. (ii) With $\bar Y=\frac1n\sum f(X_i)L(X_i)$ and
$\bar L=\frac1n\sum L(X_i)$, $\bar Y-\hat\mu^{\mathrm{sn}}_n(\bar L-1)=\hat\mu^{\mathrm{sn}}_n$
because $\hat\mu^{\mathrm{sn}}_n\bar L=\bar Y$.
\end{proof}

The correlation that matters is the one between the \emph{weighted} quantities under $\Qt$.
The linear control $S_T$ is weak for deep out-of-the-money calls even under $\Q$ (exact
$\rho^2=0.10$ at $K=150$ and $0.003$ at $K=200$, against $0.78$ at $K=100$), and after the GHS
shift it is weaker still ($\tilde\rho^2=0.003$ and $8\times10^{-7}$), because the drift removes
the linear part of $\log G$ (\cref{prop:is-optimal-drift}). A
nonlinear control that tracks the payoff on the important paths survives, such as the geometric
Asian call for the Asian (\cref{prop:as-control}). In \cref{tab:is-lab} the control dominates
near the money, importance sampling in the tail, and together they win at every strike. The
self-normalised estimator suits weights known up to a constant (\exref{ex:is-self-normalized}).

\begin{table}[ht]
  \centering
  \small\setlength{\tabcolsep}{4.5pt}
  \begin{tabular}{@{}rr@{\quad}cccc@{\quad}ccc@{}}
    \toprule
    & & \multicolumn{4}{c}{standard error, $n=2^{17}$} & \multicolumn{3}{c}{variance reduction} \\
    $K$ & price & plain & control & IS & both & control & IS & both \\
    \midrule
    $100$ & $4.4963$ & $2.0\times10^{-2}$ & $5.3\times10^{-4}$ & $6.7\times10^{-3}$ & $4.0\times10^{-4}$ & $1458$ & $9.2$ & $2560$ \\
    $120$ & $0.36113$ & $5.7\times10^{-3}$ & $3.9\times10^{-4}$ & $8.8\times10^{-4}$ & $1.4\times10^{-4}$ & $220$ & $42$ & $1621$ \\
    $140$ & $0.012755$ & $1.0\times10^{-3}$ & $1.7\times10^{-4}$ & $4.1\times10^{-5}$ & $1.5\times10^{-5}$ & $34$ & $618$ & $4359$ \\
    $160$ & $2.7333\times10^{-4}$ & $1.4\times10^{-4}$ & $4.7\times10^{-5}$ & $1.0\times10^{-6}$ & $6.6\times10^{-7}$ & $9.2$ & $1.9\times10^4$ & $4.7\times10^4$ \\
    \bottomrule
  \end{tabular}
  \caption{Arithmetic Asian call, test case, $16$ fixings: plain Monte Carlo, the geometric
    control, importance sampling with the GHS drift, and both (control fitted in sample with
    \texttt{variance.control\_variate}). Our implementation (\texttt{default\_rng},
    \texttt{SeedSequence(2026)}, \texttt{importance.asian\_laboratory}); price: combined
    estimate. Only $9$ of $2^{17}$ plain paths pay at $K=160$, so the variances of plain Monte
    Carlo and of the control alone are estimated from the importance sample
    ($\EQ[h]=\Et[hL]$, optimal coefficient).}
  \label{tab:is-lab}
\end{table}

\demo{09-importance-sampling.ipynb}{Asian drifts, effective sample sizes, this laboratory.}

\furtherreading
\Cref{fig:is-densities,fig:is-contours} and \cref{rem:is-share} reproduce the course's
slides. \textcite{glasserman-mc} covers the likelihood ratio and tilting in Section~4.6.1 and
the GHS drift in Section~4.6.2; \textcite{asmussen-glynn} treat importance sampling in
Section~V.1 and rare events in Chapter~VI (efficiency VI.1, tilting VI.2, cross-entropy VI.8).
Cross-entropy is due to \textcite{rubinstein-ce} (tutorial: \textcite{deboer-ce});
\textcite{dembo-zeitouni}, Chapter~2, proves Cramér's theorem.
